\section{Results and Discussion}

\subsection{What Works Well}
Several aspects of the system are already strong from both an embedded-systems
and a software-engineering perspective.

First, the task architecture is coherent. The project does not implement
assignment features as unrelated code fragments; instead, all major functions
fit into a single RTOS dataflow model. The sensor task is the primary producer,
the UI tasks are consumers, the TinyML task is a transformer, and the CoreIOT
task is the cloud sink. This architectural clarity is a major strength of the
system.

Second, the use of synchronization primitives is appropriate. Queues are used
for data movement, mutexes are used for genuine shared-resource protection, and
the binary semaphore is used for a one-time readiness signal. This is a better
design than trying to replace all state with semaphores or relying on broad
global-variable access.

Third, the AP dashboard is practically useful. It is not limited to a single
HTML page with one button. It offers Wi-Fi scanning, profile management, live
telemetry, LED control, brightness control, and recent-history visualization.
This gives the project a user-facing completeness that many small embedded
projects do not reach.

Fourth, the project demonstrates good feature integration. TinyML is not
isolated from the rest of the system; it affects the LCD, the dashboard, and
the cloud payload. Likewise, CoreIOT is not treated as a one-way uploader only;
it also influences local behavior through RPC.

\subsection{Tradeoffs in the Current Design}
Although the implementation is successful overall, the design includes several
explicit tradeoffs.

\subsubsection*{AP usability versus STA cloud publishing.}
The same board cannot simultaneously behave as a simple AP-only configuration
node and as a pure Internet-connected cloud device without state transitions.
The chosen design resolves this by prioritizing AP mode for local setup and STA
mode for cloud publishing. This is practical, but it introduces mode-switching
complexity and makes some functionality context-dependent.

\subsubsection*{Fresh data versus full history.}
The queues mostly use single-slot ``latest value'' semantics. This is the right
decision for UI responsiveness and low memory consumption, but it means the
system does not preserve a complete time series for every consumer. The web
task therefore implements its own history buffer for visualization.

\subsubsection*{Prediction stability versus latency.}
The TinyML task batches ten sensor samples before inference. This reduces noise
and gives the model smoother features, but it increases prediction latency to
roughly 20 seconds. For a weather-oriented embedded node, this is acceptable,
but for a faster reactive use case it would be too slow.

\subsubsection*{Reduced global state versus local module state.}
The project has reduced shared global variables significantly by introducing an
application context and protected shared state. However, some file-local static
variables still exist, especially inside the web module. This is not a design
error; it is a tradeoff between strong encapsulation and total elimination of
all persistent module state.

\subsection{Current Weaknesses}
Several limitations remain in the current system.

\subsubsection*{Timing validation is mostly analytical.}
The code is structured well for soft real-time execution, but hard timing
evidence is limited. There are no per-task runtime statistics, no deadline miss
counters, and no direct WCET measurements.

\subsubsection*{Cloud connectivity still depends on external network conditions.}
The CoreIOT task is correctly gated by STA readiness, but MQTT success still
depends on Wi-Fi quality, broker availability, and network environment. This is
expected in practice, but it prevents deterministic end-to-end guarantees.

\subsubsection*{TinyML evaluation is deployment-focused.}
The code clearly deploys and executes the model, but the repository itself does
not yet contain a complete experimental section with confusion matrices, test
sets, or repeated on-device accuracy runs. That evaluation should be added in
the final report using training and test results from the model-development
pipeline.

\subsubsection*{Time synchronization is prepared but not fully integrated.}
The time utility module exists, but active NTP synchronization is not visible in
the main execution flow. As a result, some time-derived TinyML features may rely
on fallback constants and some telemetry timestamps may fall back to RTOS tick
time.

\subsubsection*{The web chart is intentionally lightweight.}
The dashboard chart fulfills the current monitoring requirement, but it is not
yet a full monitoring frontend. It keeps only a short recent history and does
not persist records across resets or support export.

\subsection{Future Development Directions}
The current system already satisfies the assignment, but several extensions are
natural next steps:
\begin{itemize}
  \item add explicit runtime profiling for each task,
  \item integrate NTP synchronization into the STA workflow,
  \item increase the length and persistence of the web history buffer,
  \item support offline buffering and delayed CoreIOT upload,
  \item extend cloud RPC to control more local features,
  \item add more sensors or actuators using the same queue-based architecture,
  \item improve TinyML evaluation with clearer dataset and accuracy reporting.
\end{itemize}

In summary, the system is already well structured and functionally complete, but
the strongest next improvements would focus on measurement rigor, connectivity
robustness, and richer data analysis.
